\subsection{An obstruction at the curvature-matched scales}
\label{hybrid:subsection}

The coherent approximations at source-cell depth \(n\) have scalar error
\(O(2^{-2n})\). Their natural comparison with frequency \(R=2^N\) is
therefore \(N=2n\), rather than \(N=n\). The following calculation shows
that the matching profiles can obstruct both existing estimates at these
matched scales. It does not exclude an estimate coupling the two arguments
in a genuinely new way.

\begin{lemma}\label{hybrid:collision}
Let \(\mu\) be a homogeneous dyadic probability: at every integer depth
\(k\), all occupied squares have mass \(2^{-F(k)}\). Assume that the
conditional energies below are finite. Fix \(0<\gamma<1/2\), put
\(\tau=1+2\gamma\), and use the coherent quantity
\[
 Z_n=2^{-n(1+4\gamma)}
 \sum_{Q\text{ at depth }n+1}\mu(Q)I_\tau(\mu_Q).
\]
For every integer \(N\ge n+1\),
\begin{equation}\label{hybrid:collision-bound}
 Z_n\ge c_\gamma
 2^{\tau N+F(n+1)-F(N)-n(2\tau-1)}.
\end{equation}
In particular, if \(N=2n\), then
\begin{equation}\label{hybrid:midpoint}
 Z_n\ge c_\gamma\,2^{F(n)-F(2n)+n}.
\end{equation}
\end{lemma}
\begin{proof}
Fix an occupied square \(Q\) of depth \(n+1\). It contains exactly
\(2^{F(N)-F(n+1)}\) occupied depth-\(N\) descendants, each with
conditional mass \(2^{F(n+1)-F(N)}\). Two independent points under
\(\mu_Q\) lie in the same such descendant with probability
\[
 2^{F(n+1)-F(N)}.
\]
Their separation is at most \(\sqrt2\,2^{-N}\). These pairs alone
give
\[
 I_\tau(\mu_Q)\ge
 2^{-\tau/2}2^{\tau N+F(n+1)-F(N)}.
\]
Summing with the parent weights, which total one, proves
\eqref{hybrid:collision-bound}. For \(N=2n\), the terms containing
\(\tau\) reduce to \(n\); also \(F(n+1)\ge F(n)\).
This gives \eqref{hybrid:midpoint}.
\end{proof}

Suppose even depths \(N_j\) realize a template as in
Theorem~\ref{realization:theorem}:
\[
 \frac{F(N_jx)}{N_j}\longrightarrow x+g(x)
 \quad\hbox{uniformly},\qquad g(1)=a.
\]
Equation~\eqref{hybrid:midpoint} gives
\begin{equation}\label{hybrid:exponent}
 \liminf_{j\to\infty}\frac1{N_j}\log_2 Z_{N_j/2}
 \ge g(1/2)-a.
\end{equation}
Consequently, when \(g(1/2)=b/2\) and \(b>2a\), the coherent
quantities grow along this subsequence. The coherent summability
criterion cannot hold. This calculation is independent of the legal
choice of \(\gamma\). For the realized \(d\)-Frostman measures,
the required conditional energies are finite whenever
\(1+2\gamma<d=1+a\).
Integer rounding or a fixed additive change of the matching depth only
changes the exponent by \(O(1)\). A strictly positive limiting exponent
also persists under sufficiently small fixed proportional changes of
that depth.

\paragraph{Checking the matching templates.}
For the single-collapse witness of
Section~\ref{hardgap:one-witness}, with \(0<a<b\le1\), its first
affine piece ends at
\[
 P=\frac{1+a}{1+b}>\frac12.
\]
Hence \(g(1/2)=b/2\). Whenever
\[
 b\ge b_F(a)=\frac{2a}{1-2a},\qquad 0<a\le\frac14,
\]
its limiting chain cost \(C_1(a,b)\) is at least \(a\).
Since \(b_F(a)>2a\), the same witness obstructs the strict chain-cost
test and the coherent summability estimate at the matched scales.

For the two-collapse witness of
Section~\ref{hardgap:two-witness}, its initial affine piece ends at
\[
 L(a,b)=
 \frac{(1+a)(1-b)(1+2a)}
 {(1-a)(1+2b)(1+b)}.
\]
This function decreases in \(b\). At
\(b=b_{\mathrm c}(a)=(1+2a)/(4+2a)\), direct simplification gives
\[
 L(a,b_{\mathrm c}(a))-\frac12
 =\frac{8a^2+11a-1}{2(1-a)(4a+5)}.
\]
For \(\alpha=(3-\sqrt7)/4\), its numerator equals \(23\alpha-2>0\);
the last inequality follows from \(61^2>23^2\cdot7\).
The numerator increases for positive \(a\). Thus
\[
 \alpha\le a\le a_c,\qquad
 2a\le b\le b_{\mathrm c}(a)
 \quad\Longrightarrow\quad L(a,b)>\frac12.
\]
In particular \(g(1/2)=b/2\) throughout this range.
For \(\alpha<a\le a_c\) and
\(r_2(a)\le b\le b_{\mathrm c}(a)\), the proved matching cost
satisfies \(C_2(a,b)\ge a\), while \(b>2a\).
Equation~\eqref{hybrid:exponent} again makes both existing sufficient
tests fail on the same realized example, with their indices matched
by curvature.

\paragraph{The zero-exponent boundary can also be realized without decay.}
For \(0<a\le\alpha\), use the two-collapse witness with \(b=2a\).
It is in the witness's parameter domain: \(\alpha<1/8\) implies
\(4a^2+6a<1\), equivalently \(2a<b_{\mathrm c}(a)\).
It has cost at least \(a\), and
\[
 L(a,2a)-\frac12
 =\frac{1-5a}{2(1-a)(1+4a)}>0.
\]
Thus \(g(1/2)=a\). Uniform template convergence alone only gives a
zero limiting exponent in \eqref{hybrid:exponent}; it does not prove
failure of summability. The following modification of the realization
does prove that failure.

Write \(d=1+a\), \(D=1+2a<2\), and \(h(x)=x+g(x)\).
Suppose a block starts at depth \(M\) with cumulative value \(dM\).
Choose the next even endpoint \(N\) sufficiently large and, on
\([M,N]\), define
\[
 \mathcal F(t)=
 \min\{Nh(t/N),\,2t-(2-d)M\}.
\]
The template initially equals \(Dx\), on an interval extending past
\(1/2\). The two entries first meet at
\[
 t_0=\frac{2-d}{2-D}M.
\]
Choose \(N\) so large that \(t_0<N/2\) and that this meeting occurs
in the initial linear part. Before \(t_0\) the second entry is smaller.
Afterwards it remains at least the first, since it has slope two
and the template is 2-Lipschitz. Therefore
\[
 \mathcal F(N/2)=DN/2,\qquad \mathcal F(N)=dN.
\]
This block joins the preceding one continuously, is nondecreasing and
2-Lipschitz, and obeys \(dt\le\mathcal F(t)\le Dt\).
Choose successive blocks with \(M/N\to0\). Away from an initial
interval of length \(O_{a}(M)\) the block agrees exactly with its
template, so the same uniform realization limit holds.

Use the product-digit construction with integer exponent
\(F(k)=2\lfloor\mathcal F(k)/2\rfloor\), as in
Theorem~\ref{realization:theorem}. Its uniform \(d\)-Frostman and
upper-cover bounds follow from the same barriers. Its exact Hausdorff
and packing dimensions are \(d,D\), using the endpoint and midpoint
scales, respectively, in that theorem's dimension argument.
The rounding error is less than two at integer depths, and hence
\[
 F(N/2)-F(N)+N/2\ge-2.
\]
Equation~\eqref{hybrid:midpoint} implies
\(Z_{N/2}\ge c_\gamma/4\) along infinitely many blocks. Thus the
coherent summability criterion fails even at this zero-exponent
boundary. The realization still has the matching limiting chain cost
at least \(a\).

Fixed positive-mass cylinder restrictions preserve the calculation:
their later-depth exponent is \(F(k)-F(k_0)\), and the subtracted
constant cancels in \eqref{hybrid:midpoint}. The separated source
and pin restrictions supplied by the realization construction are
therefore compatible with these examples.

This result concerns the two existing estimates, not the actual distance
laws. In particular, one cannot infer a hybrid improvement merely because
the coherent comparison looks favorable at depth \(N\) where the Fourier
profile is unfavorable at frequency \(2^N\): the directly matching
coherent depth is \(N/2\), where the examples above are unfavorable.
Controlling a different bridge between the two approximation schemes
could still improve the distance theorem.
